\chapter[The Compute: RISC-V, Accelerators, and the Memory Arbitrage]{The Compute: RISC-V, Accelerators, and the Memory Arbitrage}
\chaptermark{The Compute}
\label{ch:compute}

The second element is the machine the models run on. The Western AI machine is a proprietary stack---x86/Arm hosts, NVIDIA accelerators, CUDA, NVLink, HBM---each layer a licensed chokepoint. China's counter-design, now visible end to end, is an \emph{open} machine: open instruction sets, multi-vendor domestic accelerators unified by open interconnect and open software, and a memory system that trades scarce bandwidth for abundant capacity. In June 2026 this ceased to be a projection: an all-CPU, GPU-free Chinese machine took the No.~1 position on both of the world's canonical supercomputing benchmarks. The claim this book draws from that is narrower than the headlines suggest and, properly stated, more interesting: not that the GPU is obsolete---it is not, and Section~\ref{sec:notobsolete} says so plainly---but that the CPU corner of the design continuum can reach performance parity while carrying a software and application-ecosystem advantage the accelerator corner cannot easily replicate. This chapter treats the layers in turn---the instruction set, the machine that proved the architecture, the template it establishes, the accelerator fleet, the memory arbitrage, and the quiet defection of the Western ecosystem itself to the open ISA.

\section{State-mandated architectural independence: RISC-V}
\label{sec:riscv}

China is actively organizing its domestic industry to circumvent the Western ISA monopolies---Arm and x86---that form the bedrock of Western computing. In March 2025, eight government bodies including MIIT, the CAC, and MOST jointly drafted the first national policy guidance for RISC-V adoption nationwide~\cite{riscv-policy}; the 15th Five-Year Plan (2026--2030) mandates ``extraordinary measures'' for integrated circuits and basic software self-reliance, with accelerated deployment of the open RISC-V instruction set as a core mechanism~\cite{fyp15}. Because RISC-V is an open standard, its deployment insulates Chinese silicon design from sanctions at the ISA layer permanently: there is no license to revoke.

The execution is characteristically whole-of-nation. The Chinese Academy of Sciences fields one of the world's largest RISC-V development organizations---over 600 hardware and 400 software researchers---whose open-source XiangShan cores now benchmark at Arm Neoverse-class performance (Kunminghu: 6-wide decode, four RVV~1.0 vector units, $\sim$15 SPECint/GHz), with the openRuyi OS as the first RVA23-native software platform~\cite{xiangshan,cas-riscv}. Alibaba's T-Head ships the XuanTie C930 server-class RISC-V CPU; SpacemiT's K3 pairs an eight-core RISC-V cluster with a 60-TOPS AI engine~\cite{c930,spacemit}. Within RISC-V International (Swiss-incorporated precisely to be sanction-proof), 12 of 24 premier members are Chinese~\cite{csis-riscv}. The RISC-V Summit China 2026 (Shenzhen, October 18--20) is projected to host over 450 exhibitors and 3,600 industry professionals, its call for topics explicitly focused on ISA standardization of matrix and vector extensions across the full AI computing stack~\cite{summit-2026}. Those extensions---the IME, VME, and AME task groups---will give the open ISA native matrix operations, and Chinese vendors are the largest bloc shaping them~\cite{matrix-ext}. Their schedules and roles must be stated separately, because they are commonly presented as sharing one near-term ratification horizon and they do not: IME (integrated matrix extension---matrix operations executed within the vector unit) targeted ratification in August 2026; AME (attached matrix extension---a tile engine with its own architectural state, the closest analogue to Arm's SME and the extension class the accelerated-CPU template of Section~\ref{sec:template} actually requires) carries a board-ratification target of April 2027; VME (vector--matrix interaction) follows its own track~\cite{imePlan,amePlan}. And ratification is only the first rung of a maturity ladder the strategy must climb end to end:
\begin{align*}
\text{specification} \rightarrow \text{hardware} \rightarrow{}& \text{compiler lowering} \rightarrow \text{kernel library}\\
\rightarrow{}& \text{framework integration} \rightarrow \text{fleet validation},
\end{align*}
with the AME plan itself treating simulators, ABI definition, GCC/LLVM support, and basic kernels as staged milestones and optimized software as an explicitly later problem~\cite{amePlan}. The RISC-V accelerated-CPU path is therefore real but dated: the ISA foundation completes across 2026--2027, and the software rungs above it are measured in further years---which is exactly the Fugaku-calibrated interval Section~\ref{sec:proofboundary} describes, now with its clock start pinned to primary-source schedules. Vendors such as EVAS Intelligence are already shipping full-stack RISC-V AI ``supernode'' systems---a self-developed cloud AI chip with block-quantized FP8, 3.2\,TB/s per-chip interconnect, and claimed scaling to $10^5$-card clusters~\cite{evas}. Two billion RISC-V chips have shipped; current forecasts project 36 billion units annually by 2031, with AI-accelerator SoCs alone accounting for roughly 70\% of RISC-V silicon revenue~\cite{shd-2026}.

The strategic reading: the ISA layer is being moved from the rent column to the commons column, exactly as model weights were in Chapter~\ref{ch:model}. What the ISA strategy needed, as of 2025, was an existence proof that a domestically fielded, non-x86, non-NVIDIA machine could stand at the absolute top of world computing. It arrived ahead of schedule---on Arm rather than RISC-V, which, as Section~\ref{sec:template} argues, is precisely the point.

\section{LineShine and the LX2: the GPU-free machine takes No.~1}
\label{sec:lineshine}

On June 24, 2026, at ISC in Hamburg, the 67th TOP500 list opened a new era twice over. LineShine (\emph{Ling Sheng}, ``intelligent splendor''), built by the National Supercomputing Center in Shenzhen, debuted at \textbf{No.~1 on HPL at 2.198\,EFLOP/s}---the first machine ever to sustain two exaflops of FP64, the first Chinese No.~1 since 2017, and, most consequentially for this book, \textbf{a machine containing no GPUs whatsoever}~\cite{top500-lineshine,dongarra-report}. It simultaneously took \textbf{No.~1 on HPCG} (22.00\,PF, ahead of El Capitan's 17.41 and Fugaku's 16.00)---the memory-bandwidth-bound benchmark---while posting 7.92\,EF on mixed-precision HPL-MxP, at 42.2\,MW and 52.1\,GF/W, in the same Green500 band as the US GPU flagships~\cite{top500-lineshine,dhm-long}.

The processor behind it, the Armv9 \textbf{LX2}, is the architectural pivot of this chapter: 304 cores per socket, every core carrying both SVE vector and SME matrix engines, with 32\,GB of on-package HBM at 4\,TB/s plus a large DDR5 tier---60.3\,TF FP64, 240\,TF BF16, 960\,TOPS INT8 per socket at $\sim$690\,W~\cite{dongarra-report,dhm-long}. It is Fugaku's architectural lineage (A64FX: the first general-purpose CPU with on-package HBM) carried to its logical conclusion: the accelerator's defining features---wide vectors, matrix engines, HBM, low precision---absorbed \emph{into the CPU instruction set}. Precision about what this is: the LX2 is an \emph{integrated accelerated processor}, GPU-free in product taxonomy rather than free of accelerator-class silicon. The claim it supports is correspondingly precise---not that accelerator silicon is unnecessary, but that the discrete accelerator \emph{product category}, with the separate vendor, separate memory space, separate software monopoly, and separate margin structure attached to it, is an implementation choice rather than an architectural necessity. Its designer is undisclosed; analysts read Huawei's fingerprints in the software stack, SMIC 7\,nm-class as the process candidate ``by elimination,'' and Chinese state media claim the machine runs on ``China's first domestically developed HBM''---a claim examined in Chapter~\ref{ch:semi}~\cite{lockwood,globaltimes-lineshine}.

The quantitative meaning of this machine for AI---not just simulation---is worked out in the study by Dongarra, Hoefler, and Matsuoka, ``Do We Still Need GPUs?'', in its short form~\cite{dhm-cacm} and full-length companion~\cite{dhm-long}, and the book adopts its findings as load-bearing. Four of them matter here. \emph{First}, for the traditional HPL/HPCG workload space the question is settled on the public record [V]: a CPU machine now leads both the compute-bound and the bandwidth-bound FP64 benchmarks at GPU-class energy efficiency---while transformer training and high-concurrency serving remain, as the study itself insists, projections pending direct measurement [M]. \emph{Second}, for frontier AI serving---tested at the hardest available point, a 1T-parameter MoE (Kimi-K2) at 256K context---decode reduces to a bandwidth roofline on which \emph{an LX2 socket serves the same users per TB/s as a GPU}, and prefill reduces to stated, achievable thresholds on INT8 matrix efficiency and interconnect; the all-CPU machine lands at roughly 1.3--1.4$\times$ a B200 fleet's per-stream energy under central assumptions~\cite{dhm-long}. \emph{Third}, training on the shipping part costs 4.5--8.5$\times$ a GPU fleet's energy per token at matched BF16---but every term of that deficit (the missing FP8 path, engine width, fabric ratio) is an implementation choice, not an architectural constant; an FP8-equipped successor cuts the central gap to $\sim$3$\times$, and a part with GPU-parity matrix throughput trains at parity to first order. \emph{Fourth}, and most generally: with HBM now common to both device classes, ``CPU'' versus ``GPU'' has collapsed into a single design continuum parameterized by die-area allocation between low-precision matrix throughput and high-precision generality---\emph{either corner can be provisioned to do the other's job}~\cite{dhm-cacm,dhm-long}.

Jack Dongarra's own verdict, delivered alongside his technical report on the system, states the strategic conclusion this book would otherwise have to argue: China ``can adapt to develop its own version of technology as good as---or maybe even better than---existing technology, despite US export controls''~\cite{dongarra-scmp,dongarra-report}. Export controls denied China the discrete GPU; the response was to demonstrate, at No.~1 in the world, that the functional boundary between CPU and accelerator is an implementation choice---and then to choose an implementation with no controllable product category in it.

\subsection{The proof boundary: what this machine establishes, and what it does not}
\label{sec:proofboundary}

That last sentence is the strongest one this book will make about LineShine, and it is worth stating exactly where its support ends---because the argument survives a much narrower reading than its enthusiasts give it, and is more useful stated narrowly.

\begin{figure}[htbp]
\centering
\begin{tikzpicture}[
  every node/.style={font=\scriptsize},
  band/.style={rounded corners=3pt, inner sep=5pt, align=left, text width=44mm,
               minimum height=41mm, anchor=north west},
  item/.style={font=\scriptsize\color{inkSec}},
]
\node[band, fill=sOne!12, draw=sOne, line width=0.7pt] (proven) at (0,0)
  {\textbf{\normalsize Measured [V]}\\[3pt]
   HPL 2.198\,EF (No.\,1)\\
   HPCG 22.00\,PF (No.\,1)\\
   HPL-MxP 7.92\,EF\\
   42.22\,MW, 52.1\,GF/W\\
   304-core LX2, no discrete GPU\\
   A64FX lineage: STREAM parity\\
   with contemporary GPU (2018)\\
   450\,GB/s measured per domain};
\node[band, fill=sThree!12, draw=sThree, line width=0.7pt] (inferred) at (4.95,0)
  {\textbf{\normalsize Roofline-inferred [V/M]}\\[3pt]
   decode is bandwidth-bound:\\
   users served per TB/s is a\\
   \emph{ratio}, not a fitted model\\[3pt]
   $\Rightarrow$ one LX2 socket serves the\\
   same decode streams per TB/s\\
   as a GPU, at half the\\
   per-package bandwidth\\[3pt]
   only the efficiency coefficient is\\
   modelled, and swept};
\node[band, fill=sTwo!12, draw=sTwo, line width=0.7pt] (modelled) at (9.90,0)
  {\textbf{\normalsize Modelled [M]}\\[3pt]
   prefill above stated thresholds\\
   ($e_{\text{INT8}}\gtrsim0.38$, $e_{\mathrm{INT8}}\gtrsim0.90$)\\
   all-reduce break-evens 9--38\,GB/s\\
   fleet energy $\approx$1.3--1.4$\times$ B200\\
   training deficit 4.5--8.5$\times$ at\\
   matched BF16; $\sim$3$\times$ with FP8\\[3pt]
   \emph{not yet measured on the machine}};

\draw[-{Stealth[length=5pt]}, line width=1pt, color=inkSec]
  (0.2,-4.75) -- (14.3,-4.75)
  node[midway, above=0.5mm, font=\scriptsize\color{inkSec}]
  {increasing dependence on assumptions};

\node[font=\scriptsize\color{inkSec}, align=left, text width=146mm, anchor=north west] at (0,-5.15)
  {\textbf{Stage-1 claim (established):} discrete-GPU packaging is not required for world-leading HPL and HPCG
   performance at competitive energy. \quad
   \textbf{Stage-2 claim (argued, not yet measured):} the same integration extends to AI serving, and---with
   different matrix and memory provisioning---to training. \\[2pt]
   The book's \emph{strategic} claim requires only stage~1 plus the die-area allocation argument, which is
   independent of this machine.};
\end{tikzpicture}
\caption[The LineShine proof boundary]{What LineShine proves, and what it does not. The strongest form of the
architectural argument is not that a CPU has been shown to train transformers---it has not---but that the
functional boundary between CPU and accelerator has become an allocation choice, of which LineShine is the first
at-scale demonstration in the direction the incumbent does not occupy. Stated in the two-stage form this book adopts; measured values from the June 2026 TOP500 submission and the accompanying
technical report~\cite{top500-lineshine,dongarra-report}, modelled values from the companion
analysis~\cite{dhm-long}.}
\label{fig:proofboundary}
\end{figure}


The disciplined formulation is two-stage. \emph{Stage one, established}: discrete-GPU packaging is not required for world-leading HPL and HPCG performance at competitive energy efficiency [V]. \emph{Stage two, argued but not yet measured}: the same integration extends to AI serving, and---with different matrix and memory provisioning---to training [M]. LineShine does \emph{not} yet validate large-transformer training efficiency, expert-routing throughput and tail latency, distributed optimizer scaling, collective performance under AI communication patterns, kernel coverage and software productivity, serving throughput at production concurrency, or reliability and checkpoint/restart behaviour across multi-week runs. HPL and HPCG are not transformer benchmarks, and no amount of eminence in the rankings makes them so. Figure~\ref{fig:proofboundary} draws the line.

Having conceded that, three observations push back on the deflationary reading, and together they are why this book still treats the machine as a hinge rather than a curiosity.

\emph{First, the inference from stage one to the serving half of stage two is shorter than ``modelled'' suggests.} Autoregressive decode over a sparse mixture-of-experts is memory-bound by an order of magnitude---its arithmetic intensity sits far below any current socket's compute-to-bandwidth ridge---so the number of concurrent streams a socket serves is, to first order, a \emph{ratio} of delivered bandwidth to per-user KV traffic, not a fitted performance model. What is genuinely modelled is a single efficiency coefficient---written $\eta$ throughout this book and in the claim register, and defined as the \emph{fraction of a socket's peak memory bandwidth that is actually delivered to the serving workload}, so that achieved bandwidth is $\eta$ times nominal. The companion analysis sweeps $\eta$ over 0.53--0.70 rather than asserting a value, and every serving conclusion in this chapter is reported across that band~\cite{dhm-long}. The prerequisite that a CPU can \emph{drive} on-package HBM rather than merely attach it is demonstrated on deployed hardware in this exact lineage---STREAM parity with the contemporary GPU generation on A64FX in 2018, 450\,GB/s measured per memory domain on the LX2, and a world-record HPCG that is precisely a delivered-bandwidth benchmark. Prefill, the compute-bound half, is reduced not to an assertion but to \emph{stated numerical thresholds}---sustained INT8 efficiency above roughly 0.38, INT8-eligible fraction above 0.90, achieved all-reduce above 9--38\,GB/s---which is the form a falsifiable claim takes.

\emph{Second, the software gap is a measured historical quantity rather than an unknown.} The same lineage supplies the calibration: Fugaku led the benchmark lists from 2020, yet a performant AI stack for its architecture arrived years later. The correct prior for an architecture whose hardware prerequisites are demonstrated and whose software ecosystem is young is therefore not ``unproven'' but ``capability present, stack lagging by a measurable interval''---which is exactly how the West should read the LX2, and exactly the interval during which a determined national programme closes it.

\emph{Third, and most importantly, the strategic claim does not rest on this machine at all.} The load-bearing argument of Section~\ref{sec:template} is the die-area allocation argument: once HBM is common to both device classes, what distinguishes a future CPU from a future GPU is how compute silicon is divided between low-precision matrix throughput and high-precision generality, with injection bandwidth provisioned in ratio. That argument is architectural, not empirical, and it is corroborated from three independent directions at once. LineShine approaches the continuum from the CPU corner. NVIDIA's own product line approaches it from the GPU corner---Grace CPUs, LPDDR-based memory subsystems and SOCAMM modules, CUDA ported to RISC-V hosts (Section~\ref{sec:west-riscv}). And the hyperscalers' custom accelerators approach it from a third direction entirely: systolic and dataflow designs, none of them SIMT, converging on similar allocations without sharing a device category. Three actors with different constraints, different customers, and different starting points are converging on one design space. \emph{A category that three independent parties are dissolving is not a category.}

There is also a fairness point worth naming, because it recurs whenever a challenger architecture appears. The demand that LineShine prove transformer-training competitiveness before the architectural argument counts applies a standard the incumbent never met: GPUs were adopted for deep learning because they were available, programmable, and fast enough, years before anyone demonstrated they were the optimal allocation---and their FP64 capability has since been substantially removed in the tensor-heavy parts, which is itself an admission that the allocation is a choice. The correct question is not whether the CPU corner has proven itself against a standard the GPU corner was never held to; it is which corner the workload of the next decade rewards, and Chapters~\ref{ch:prize} and~\ref{ch:sizing} argue that long-context agentic scientific orchestration rewards the capacity-rich, generality-preserving corner more than the dense-FLOP one.

\subsection{The AI proof suite: what would settle it}
\label{sec:proofsuite}

Because the second stage is contested, this book states what would resolve it. Table~\ref{tab:proofsuite} specifies a measurement programme for LineShine or any LX2-class successor. It is deliberately constructed so that each row can be run by the machine's operators and published without disclosing anything about the processor's provenance, and so that a failure on any row is a genuine falsification of the corresponding claim in this book.

\begin{table}[htbp]
  \centering
  \small
  \caption{Proposed AI proof suite for an LX2-class integrated accelerated processor. Each row states the claim at risk and the outcome that would falsify it.}
  \label{tab:proofsuite}
  \begin{tabular}{p{40mm}p{50mm}p{52mm}}
    \toprule
    Measurement & Claim at risk & Falsified if \\
    \midrule
    MLPerf Inference (closed division), 1T-class MoE & serving parity per TB/s & delivered streams per TB/s fall materially below GPU submissions at matched precision \\
    \addlinespace
    Long-context decode at 256K, full KV traffic & the bandwidth-roofline argument & achieved fraction of peak collapses under the serving placement (aggregate $\ne$ per-domain) \\
    \addlinespace
    Prefill at production concurrency & the stated INT8 thresholds & sustained $e_{\text{INT8}}$ stays below 0.38 with a mature kernel library \\
    \addlinespace
    All-to-all expert routing at scale & the fabric-in-ratio discipline & achieved collectives fall below the stated break-evens on the declared layout \\
    \addlinespace
    MLPerf Training, and a multi-week frontier-class run & the time-indexed training verdict & energy per trained token stays worse than $\sim$4$\times$ a contemporary GPU fleet after an FP8-capable generation \\
    \addlinespace
    Checkpoint/restart and failure-domain behaviour at full state & multi-week operability & mean time between interruptions times restart cost makes long runs uneconomic \\
    \addlinespace
    Energy per \emph{accepted} training token, and per solved agentic task & the whole economic case & the delivered-work comparison reverses the roofline comparison \\
    \bottomrule
  \end{tabular}
\end{table}

A protocol note binds the whole table: every row should be run on both a dense 70B-class model \emph{and} a 1T-class sparse model, at short and 256K context, across batch and sequence-concurrency sweeps, reporting median and tail latency, achieved bandwidth \emph{by memory domain} (aggregate figures hide placement failures), energy measured at the system boundary rather than the socket, and the exact parallel layout and software versions---because Eq.~\eqref{eq:btoken} makes every headline number conditional on precisely those choices.

Until those numbers exist, the honest position is the one printed in Figure~\ref{fig:proofboundary}: chokepoint decay is demonstrated in its strongest observed form for the HPL/HPCG workload space---the chokepoint not breached but \emph{defined out of the machine}---and argued, with stated thresholds and a stated experimental programme, for the AI workload space that matters more.

\section{The template: from LX2 to a RISC-V accelerated CPU, in two variants}
\label{sec:template}

This section presents an author-modeled architecture and strategy analysis [M/R/S], not a measured system or an announced product. The LX2 settles feasibility of the class; what generalizes is the \emph{template}. The D--H--M analysis reduces the accelerated CPU to a short design sequence~\cite{dhm-long}: the memory system is chosen first, and its bandwidth fixes the (moderate) count of strong general-purpose cores needed to saturate it; scalable vector/matrix extensions then set per-core matrix throughput independently of core count---engine width becomes \emph{implementation-selectable within physical and system constraints}, because the ISA no longer fixes it (``selectable,'' not ``arbitrary'': width remains bounded by register and tile-state bandwidth, issue capacity, die area, power, timing closure, delivered memory bandwidth, context-switch cost, thermals, and interconnect injection---a qualification the term must always carry); interconnect injection is provisioned in ratio to the matrix rate; and one address space unifies the whole. Nothing in that sequence is Arm-specific. SVE and SME are simply the first shipping instances of scalable vector and matrix ISAs; RISC-V's ratified RVV~1.0 is the second scalable vector instance, and the AME extension now being standardized---with Chinese vendors as the largest bloc in the room, IME targeted for August 2026 and AME for April 2027 (Section~\ref{sec:riscv})---is the open ISA's SME equivalent. XiangShan-class open cores already carry four RVV units per core; an SME-class outer-product engine attached to such a core is an integration exercise of exactly the kind the LX2 has now de-risked, not an invention. For China the substitution is strictly preferable: the LX2 still sits on an Arm architectural license---a Western chokepoint in principle, however grandfathered in practice---whereas its RISC-V successor sits on nothing revocable at all. This is why the LX2's arrival on Arm is not a detour from the RISC-V strategy but its proof-of-concept: the national ISA plan (Section~\ref{sec:riscv}) supplies the destination, and LineShine supplies the demonstrated microarchitecture to port to it.

The template's second degree of freedom is the memory system, and here the two halves of this chapter join. Because the design sequence starts from the memory tier, one tape-out family yields two machines:

\begin{description}
  \item[The inference-centric variant] keeps the LX2's own memory choice---on-package HBM at GPU-generation bandwidth---and provisions the matrix engines to the serving thresholds the D--H--M study states (sustained INT8 efficiency $\gtrsim$0.38, achieved all-reduce cleared in ratio)~\cite{dhm-long}. This is the machine for the token factories of Chapter~\ref{ch:model}: decode at parity per TB/s with the GPU it replaces, prefill hidden inside the decode tier, FP64 retained for the converged AI-plus-simulation workload. Its binding input is HBM supply---the subject of Chapter~\ref{ch:semi}.
  \item[The training-centric variant] substitutes the HBM tier with massive LPDDR6X capacity---the arbitrage of Section~\ref{sec:lpddr6}---accepting the bandwidth knee in exchange for the 2.3\,TB/socket that eliminates activation recomputation, dissolves the sub-256-node capacity wall, and sources its memory from the one tier where a domestic supplier (CXMT) sits at the world's leading edge rather than four years behind it. Provision the same SME-class engines with an FP8 path---the single highest-leverage addition the D--H--M training analysis identifies~\cite{dhm-long}---and the pretraining energy deficit falls from categorical to a low single-digit factor, paid for in memory that costs a fraction of HBM per bit and is not export-controllable.
\end{description}

Same cores, same engines, same fabric discipline, same software stack; only the memory tier and the compute-to-bandwidth ratio differ. The dual-variant template is the hardware mirror of the model strategy: the HBM variant serves what exists, the LPDDR6X variant trains what is next, and both are buildable on open IP, on domestic 7\,nm-class logic, with domestic memory. A national program that executes it holds, in one product family, the replacement for the NVIDIA inference fleet \emph{and} the decentralized pretraining capability of Section~\ref{sec:lpddr6}---with no Western license anywhere in the bill of materials.

\section{The accelerator fleet and the unified ecosystem}
\label{sec:accel}

\subsection{Expelling NVIDIA, adopting the fleet}

The accelerator story of 2023--2026 is a pincer: US export controls squeezed NVIDIA's China products from above while Beijing squeezed demand for them from below. The A800/H800 were banned in October 2023; the H20 in April 2025 (a \$4.5B write-off), un-banned in August 2025 in exchange for an unprecedented 15\% revenue share to the US Treasury---whereupon the CAC summoned NVIDIA over alleged backdoors, banned major platforms from buying, and regulators ordered that new state-funded data centers use domestic AI chips exclusively, with provinces cutting electricity tariffs up to 50\% for domestic-chip facilities~\cite{nvidia-timeline,cac-ban,datacenter-mandate}. NVIDIA's China data-center revenue went to approximately zero in its Q2 FY2026; Jensen Huang's own scoreboard: ``from 95\% market share to near zero,'' followed by his November 2025 statement that ``China is going to win the AI race,'' walked back within hours~\cite{nvidia-zero,huang-quotes}.

Into the vacuum, the fleet. Huawei's Ascend line shipped $\sim$507,000 units in 2024 and $\sim$805,000 in 2025 ($\sim$650k of them 910C-class), with a disclosed three-year public roadmap---950PR/950DT in 2026 with Huawei \emph{self-developed} HBM (HiBL~1.0, HiZQ~2.0), 960 in 2027, 970 in 2028~\cite{semianalysis-ascend,huawei-roadmap}. Cambricon's revenue grew 453\% in 2025 to its first annual profit, targeting 500,000 accelerators in 2026; Moore Threads' Shanghai IPO popped $\sim$400\% on debut; Biren listed in Hong Kong; Baidu lit a 30,000-chip Kunlun P800 cluster; Alibaba's T-Head PPU matched the H20 closely enough that China Unicom deployed 16,384 of them in a single facility, and a 10,000-chip cluster of Alibaba's ``Zhenwu'' followed in April 2026~\cite{cambricon,ipos-chips,kunlun,ppu}. In May 2026, for the first time, nine domestic AI chips from seven vendors entered the ``secure and reliable'' government procurement catalogue~\cite{anke}. Domestic chips took roughly 41\% of Chinese AI-chip shipments in 2025 ($\sim$1.65M units) and are projected to pass 60\% in 2026~\cite{domestic-share}. China's aggregate intelligent-computing capacity reached 2,185 EFLOPS by mid-2026 at 71\% utilization~\cite{miit-eflops}.

\subsection{The unified interconnection ecosystem}

The fleet's weakness---fragmentation across a dozen incompatible vendors---is being engineered away by mandate and by open specification. MIIT policy directs the creation of a unified AI-chip computing interconnection ecosystem, so that domestic accelerators from competing vendors interoperate within shared clusters (the Computing Power Interconnection action plan; the ``AI+'' implementation documents)~\cite{miit-interconnect}. Huawei supplied the concrete substrate by \emph{open-sourcing} the crown jewels: the UnifiedBus~2.0 interconnect specification, released openly so that ``10,000+ NPUs work as one machine,'' with 300+ Atlas~900 SuperPoDs already deployed on UB~1.0; the full CANN software stack (its CUDA analogue, seven years in development), open-sourced in August 2025; MindSpore; and the complete SuperPoD reference architecture~\cite{huawei-ub,huawei-cann}. The announced Atlas~950 SuperPoD couples 8,192 Ascend~950DTs into 8 FP8-EFLOPS with 16\,PB/s of interconnect; the 960-generation SuperCluster targets over one million NPUs~\cite{huawei-roadmap}. One must apply the usual discount to vendor roadmaps---but the \emph{structure} is the point: where the West's answer to NVLink is a slow-moving industry consortium (UALink) that excludes China, China's answer is an open specification with state-mandated adoption. This top-down standardization actively dismantles the need for proprietary software stacks like CUDA, enforcing interoperability across the domestic supply chain---the Top Runner program, replayed at the cluster level. LineShine's own LingQi fabric---1.6\,Tb/s injection per node, $\ge$3.5\,Pb/s bisection---demonstrates the discipline at full machine scale~\cite{dongarra-report}.

\section{What the argument is not: the GPU is not obsolete}
\label{sec:notobsolete}

Two claims must be kept apart, and the enthusiasm around this architecture routinely blurs them. The first---\emph{the architectural continuum is real}---is strongly supported: by LineShine, by Grace, by TPUs and dataflow accelerators, by integrated HBM processors, and by the progressive redistribution of die area among scalar, vector, matrix, memory, and communication functions across every vendor's roadmap. The second---\emph{that the CPU endpoint will become economically optimal for frontier AI}---is unproven, and depends on software maturity, sustained matrix efficiency, interconnect, manufacturing economics, deployment reliability, and total cost of ownership. The disciplined statement is this:

\begin{quote}
LineShine demonstrates that the traditional CPU--GPU category boundary is no longer fundamental. It does not yet demonstrate that any particular point on the resulting design continuum is economically optimal for frontier AI.
\end{quote}

Nothing in this chapter should be read as predicting the obsolescence of discrete accelerators. The GPU corner is not a legacy position: it holds the deepest kernel libraries in computing, seven years of compounded framework integration, tuned collectives, the largest installed base of trained engineers, and---not least---an operational track record at fleet scale that no CPU-side machine has yet attempted for a multi-week frontier run. A book that predicted the disappearance of that would be making the same category error it accuses others of making about Chinese industrial capability.

The interesting claim is different, and it is positive rather than eliminative: \emph{the CPU corner can reach performance parity on the workloads that matter, and if it does, it arrives carrying an ecosystem advantage that the accelerator corner does not have.}

\subsection{Breadth against depth}
\label{sec:breadth-depth}

The two corners have opposite software strengths, and naming them correctly is more useful than ranking them. But first, a framing discipline the comparison itself demands: the practical alternatives are not ``a CPU'' and ``a GPU.'' They are an \emph{integrated accelerated-CPU system} and a \emph{heterogeneous CPU--GPU system} with coherent or semi-coherent memory and a decade of orchestration maturity---and the heterogeneous system does not need to acquire the application universe, because it can keep irregular control, I/O, meshing, and legacy code on its host while accelerating selected kernels. The honest comparison is therefore between complete workflows:
\begin{equation}
T_{\mathrm{workflow}} \;=\; T_{\mathrm{host}} + T_{\mathrm{accelerator}} + T_{\mathrm{movement}} + T_{\mathrm{synchronization}} + T_{\mathrm{orchestration}},
\label{eq:tworkflow}
\end{equation}
where every term is wall-clock time for one complete unit of work: $T_{\mathrm{host}}$ is time spent on the general-purpose processor (I/O, meshing, control flow, sparse assembly, the irregular tail); $T_{\mathrm{accelerator}}$ is time in the accelerated kernels themselves; $T_{\mathrm{movement}}$ is data transfer between memory spaces; $T_{\mathrm{synchronization}}$ is waiting at barriers and collectives; and $T_{\mathrm{orchestration}}$ is scheduling, launch, and runtime overhead. The terms are additive only to the extent they are not overlapped---a well-tuned heterogeneous system hides part of $T_{\mathrm{movement}}$ behind $T_{\mathrm{accelerator}}$, and the equation should be read as an upper bound that good engineering erodes rather than as an identity.
The integrated design wins exactly where the removal of the movement, synchronization, and orchestration terms outweighs the discrete accelerator's kernel-depth advantage on $T_{\mathrm{accelerator}}$. For kernel-dominated dense training, the last term structure favors the discrete accelerator and will for years. For the interleaved simulation--analysis--inference workflows of agentic science (Chapter~\ref{ch:prize}), the middle terms are a large and growing fraction of wall-clock, and they are the terms the integrated design deletes. Everything below should be read against Eq.~\eqref{eq:tworkflow}, not against a device-category comparison.

A second discipline: ``runs natively'' is not one property but four, and the breadth advantage is strongest on the first two:
\begin{description}\itemsep2pt
  \item[Binary/ABI portability] --- existing binaries and build systems execute unchanged. The integrated CPU wins this outright; the heterogeneous system wins it too, on its host side.
  \item[Source portability] --- code recompiles without restructuring. The integrated corner wins broadly; accelerator offload requires directive or kernel-language surgery.
  \item[Performance portability] --- the recompiled code actually exploits the matrix engines. \emph{Not} automatic on either corner: legacy scientific code may compile on an accelerated CPU, while exploiting SME/AME-class engines still requires library substitution, data-layout changes, mixed-precision reformulation, compiler maturity, and sometimes algorithm redesign. This is the rung where the maturity ladder of Section~\ref{sec:riscv} bites both sides.
  \item[Operational portability] --- schedulers, monitoring, checkpointing, failure handling, and operator skills transfer. The CPU corner's genuine decades-deep advantage; conversely, mature framework abstractions preserve much application-level investment on the GPU side even when the vendor retunes internal kernels.
\end{description}

The accelerator corner's advantage is \emph{depth}: an extraordinarily well-optimized vertical stack for a narrow set of operations---dense and block-sparse matrix multiply, attention, collectives, the reduction and normalization kernels around them---refined over a decade with enormous investment, integrated into every framework, and backed by profiling and debugging tools built for that stack specifically. On the kernels that matter for a transformer, this is a genuine and durable moat, and the LineShine class is years behind it (Section~\ref{sec:proofboundary}).

The CPU corner's advantage is \emph{breadth}, and it is under-priced because it is invisible in benchmark tables:

\begin{description}\itemsep2pt
  \item[One address space, no host--device split.] There is no explicit data movement, no separate allocator, no pinned-memory dance, no kernel-launch boundary, and no class of bug that exists only because two memories must be kept coherent by hand. For workflows that interleave simulation, analysis, and inference---which is what agentic science actually is (Chapter~\ref{ch:prize})---this removes the dominant source of both engineering effort and lost performance.
  \item[The existing software universe runs natively.] Fifty years of accumulated scientific, numerical, systems, and enterprise code executes without porting. In a typical scientific application the accelerated kernels are a small fraction of the source; the rest is I/O, meshing, control flow, sparse assembly, string handling, and glue---and on the accelerator corner that remainder either runs on the host anyway or must be rewritten. The CPU corner pays no porting tax on the 90\% of code nobody writes papers about.
  \item[The irregular tail is native.] Sparse operations, indirect addressing, branch-heavy control flow, dynamic data structures, mixed-precision solvers, graph traversal, and irregular communication are what much of computational science actually consists of, and they are exactly where accelerators are weakest. As AI itself turns sparser---mixture-of-experts routing, dynamic depth, retrieval, tool dispatch---this stops being a scientific-computing footnote and becomes an AI property.
  \item[Mature operations.] Operating systems, schedulers, containers, filesystems, debuggers, profilers, checkpointing, RAS, and fleet management are decades matured on the CPU corner. Multi-week reliability is one of the seven open questions in Table~\ref{tab:proofsuite}, and it is the question the CPU corner is best equipped to answer.
  \item[No kernel-language lock-in.] Code written to a scalable vector and matrix ISA is portable across vendors implementing that ISA---which is the entire point of Section~\ref{sec:riscv}, and the layer at which the rent of Chapter~\ref{ch:theory}'s Table~\ref{tab:rent} currently sits.
  \item[One fleet instead of two.] A converged machine serves simulation, AI training, AI serving, and data analytics from one procurement, one operations team, and one power envelope. For a national laboratory or a mid-size country, the elimination of the two-fleet problem is worth more than a percentage of peak.
\end{description}

The strategic bet the CPU corner represents is therefore a bet about \emph{which kind of software capital compounds faster}. Depth is a body of kernels that must be substantially rewritten every architecture generation anyway---and the industry rewrites it, repeatedly, whenever precision formats, attention mechanisms, or sparsity patterns change. Breadth is the accumulated application capital of half a century, and it does not have to be rewritten at all if the matrix engines arrive inside the ISA. If matrix throughput becomes an implementation choice rather than a device category---which is the argument of Section~\ref{sec:template}---then the corner that already holds the application universe has to acquire a decade of kernel depth, while the corner that holds the kernel depth has to acquire the application universe (or, per Eq.~\eqref{eq:tworkflow}, keep hosting it and pay the movement and orchestration terms forever). Both are hard. This book's judgement, stated as a judgement [S], is that the first is the more tractable problem, because kernels are a finite and well-specified target with a large and motivated workforce, whereas the application universe is neither finite nor specified.

The finite-kernel claim itself deserves an immediate qualification: the kernel set is finite at any moment but the target moves---across precision modes, sparsity patterns, expert-routing schemes, dynamic shapes, multimodal operators, graph compilation, distributed layouts, optimizer variants, numerical-reproducibility requirements, and failure recovery. ``Finite'' therefore means \emph{bounded and specifiable per generation}, not fixed. The empirically decidable form of the bet is: \emph{which side reduces the total lifecycle cost of supporting a changing workload portfolio faster?} That question is now a registered forecast in Appendix~\ref{app:dashboard}, with observable indicators (time-to-competitive-kernel for each new attention or sparsity mechanism on each corner; fraction of new model architectures served within a fixed window (say six months) of release; operator-hours per delivered petaflop-month), rather than a rhetorical flourish.

That is also why the strategic reading of LineShine survives the deflationary correction. The machine's significance is not that it retires the GPU. It is that a serious industrial actor has now demonstrated, at the top of the world rankings, that the ecosystem-rich corner of the continuum can be built at exascale---and if that corner also reaches AI parity, the competition between corners will be decided not by peak FLOPs but by which one is easier to program, operate, and converge. Those are the axes on which the CPU corner has never been behind.

\section{Architecture theory: what to report, and why}
\label{sec:archtheory}

Because this chapter argues from architecture rather than from benchmarks, it owes the reader the analytical vocabulary it is using. Every claim here reduces to the roofline framework: performance is bounded by $\min(\text{peak compute},\ \text{operational intensity}\times\text{achieved bandwidth})$, and the workload's operational intensity---FLOPs per byte of memory traffic---determines which bound binds~\cite{roofline2009}. Autoregressive decode over a sparse MoE has an operational intensity an order of magnitude below any current socket's ridge point, which is not a modelling choice but an arithmetic property of reading weights once per generated token; that is why decode reduces to bandwidth and why the per-TB/s comparison of Section~\ref{sec:proofboundary} is a ratio rather than a fit. Prefill sits above the ridge, which is why it reduces to matrix-engine thresholds instead. The memory-wall literature explains why the ridge point has drifted upward for thirty years and why capacity-first designs are a rational response; the communication-avoiding-algorithms literature explains why the fabric-in-ratio discipline of Section~\ref{sec:template} is a first-order requirement rather than a detail; and energy-roofline analysis explains why the right unit for the comparisons of Chapter~\ref{ch:energy} is joules per delivered token rather than watts per socket.

One refinement makes the decode-parity claim of Section~\ref{sec:proofboundary} properly conditional. ``Streams per TB/s'' is comparable across devices only under matched conditions, because the per-token byte traffic is itself a sum of terms that batching, caching, and placement all move:
\begin{equation}
B_{\mathrm{token}} \;=\; \frac{B_{\mathrm{active\ weights}}}{N_{\mathrm{effective\ batch}}} + B_{\mathrm{KV\ read}} + B_{\mathrm{KV\ write}} + B_{\mathrm{activations}} + B_{\mathrm{routing}} + B_{\mathrm{collective}},
\label{eq:btoken}
\end{equation}
Each term is bytes moved per generated token: $B_{\mathrm{active\ weights}}$ is the weight traffic for the experts actually activated, divided by $N_{\mathrm{effective\ batch}}$ (the number of concurrent sequences over which one weight read is amortized---``effective'' because expert routing means different sequences in a batch may need different experts, so the achieved amortization is below the nominal batch size); $B_{\mathrm{KV\ read}}$ and $B_{\mathrm{KV\ write}}$ are traffic to and from the key--value cache holding the attention state of the context so far; $B_{\mathrm{activations}}$ is intermediate tensor traffic; $B_{\mathrm{routing}}$ is the metadata moved to dispatch tokens to experts; and $B_{\mathrm{collective}}$ is the per-token share of cross-device communication under the declared parallelism layout. Throughput is then bounded on both sides:
\begin{equation}
\mathrm{tokens/s} \;\le\; \min\!\left(\frac{F_{\mathrm{achieved}}}{F_{\mathrm{token}}},\; \frac{\mathrm{BW}_{\mathrm{achieved}}}{B_{\mathrm{token}}}\right),
\label{eq:tokbound}
\end{equation}
where $F_{\mathrm{achieved}}$ is the arithmetic rate the device actually sustains on this workload (FLOP/s, not peak), $F_{\mathrm{token}}$ is the arithmetic work required per generated token, and $\mathrm{BW}_{\mathrm{achieved}}$ is delivered memory bandwidth under the workload's real access pattern (bytes/s, again not peak). The left-hand ratio is the compute bound and the right-hand ratio the bandwidth bound; whichever is smaller is what the machine delivers. At low batch on a sparse MoE the amortized active-weight term dominates $B_{\mathrm{token}}$, pushing the system onto the bandwidth bound; at extreme context, KV traffic dominates instead; and MLA, GQA, prefix sharing, expert caching, and quantization each move the boundary by shrinking one term or another. The per-TB/s parity of Section~\ref{sec:proofboundary} is therefore asserted only under matched active bytes per token, effective precision, cache behavior, batching, expert residency, KV representation, context length, routing, and collective pattern---which is exactly what the proof suite's protocol requires disclosed, and why MLPerf alone (which may not exercise the long-context sparse regime that carries the strategic claim) does not settle it.

Accordingly, for every design point this book proposes or evaluates, the quantities that should be reported---and which Appendix~\ref{app:lpddr} now carries or names as owed---are: delivered bandwidth under representative access patterns rather than peak; all-to-all and collective injection requirements at the declared parallelism layout; power per \emph{delivered} token; performance under expert imbalance, which is the sparse workload's most under-reported failure mode; checkpoint and restart overhead at full state; utilization loss from failure and maintenance; and software maturity and operator labour, which are costs even when they are not line items. A design compared on peak FLOP/s and nominal bandwidth alone has not been compared.

\section{The software and operations moat, made measurable}
\label{sec:lifecycle}

The breadth-versus-depth bet of Section~\ref{sec:breadth-depth} was registered as a forecast about lifecycle cost; that cost must therefore stop being a qualitative residual and become an object with a definition. Here it is:
\begin{align}
C_{\mathrm{lifecycle}} \;={}& C_{\mathrm{hardware}} + C_{\mathrm{energy}} + C_{\mathrm{porting}} + C_{\mathrm{optimization}} \nonumber\\
&+ C_{\mathrm{regression}} + C_{\mathrm{operations}} + C_{\mathrm{downtime}} + C_{\mathrm{migration}} + C_{\mathrm{decommission}},
\label{eq:lifecycle}
\end{align}
where $C_{\mathrm{hardware}}$ and $C_{\mathrm{energy}}$ are the acquisition and electricity costs that procurement compares, $C_{\mathrm{porting}}$ is the one-time cost of making the workload run at all, $C_{\mathrm{optimization}}$ the cost of making it run well, $C_{\mathrm{regression}}$ the recurring cost of keeping it running well across toolchain releases, $C_{\mathrm{operations}}$ the staffing of the fleet, $C_{\mathrm{downtime}}$ the value of lost availability, $C_{\mathrm{migration}}$ the cost of moving the workload off the platform at end of life, and $C_{\mathrm{decommission}}$ the disposal and data-transfer cost of retiring it. The first two are what procurement compares; the remaining seven are where platform advantage actually lives. CUDA's moat, decomposed, is a claim about the seven: lower porting cost (everything already runs), lower optimization cost (kernels exist), lower regression burden (the vendor absorbs it each release), mature operations tooling, less downtime, and---the term nobody prices until they pay it---migration cost \emph{away}, which is the lock-in rent collected at the end of the relationship rather than during it.

For each candidate architecture, the measurable schedule that would let a buyer compare Eq.~\eqref{eq:lifecycle} across platforms: time to first correct execution of a representative workload; time to 50\%, 75\%, and 90\% of attainable performance (the performance-portability curve, which distinguishes ``compiles'' from ``performs''---the four portability types of Section~\ref{sec:breadth-depth} made quantitative); source lines modified and architecture-specific code fraction; compiler and library defects encountered per port; person-months of optimization; regression burden across toolchain releases; operator staffing per thousand nodes; unplanned-interruption rate; checkpoint/restart loss; performance portability across hardware generations; and the cost of migrating one production workload off the platform. The comparative cases that would anchor the table span the industry's whole experience: NVIDIA/CUDA (the incumbent baseline), AMD/ROCm (the best-documented catch-up attempt), Huawei Ascend/CANN/MindSpore (the open-sourced challenger), A64FX/Fugaku (the direct historical calibration for the LX2 class---and the source of this book's ``capability present, stack lagging by a measurable interval'' prior), OpenXLA and the compiler-abstraction bet, the RISC-V vector/matrix toolchains as they mature through the Section~\ref{sec:riscv} ladder, and an LX2-class machine itself when access permits.

Two institutional observations complete the section. MLPerf provides the time-to-train anchor---and its v6.0 round now includes a DeepSeek-V3 MoE pretraining benchmark, making it directly relevant to the sparse workloads this book cares about~\cite{mlperf-v6}---but MLPerf measures the \emph{result} of software maturity, not its cost: a heroic two-hundred-engineer port scores identically to a recompile. The complementary instrument the field lacks, and which this book now formally proposes as a deliverable of the Chapter~\ref{ch:consortium} consortium, is an \emph{AI-HPC lifecycle benchmark}: a standardized workload portfolio whose reported metrics are the schedule above---engineering effort, defect counts, regression burden, operational staffing---audited across platforms on a fixed cadence. It would do for the software moat what TOP500 did for FLOPs: convert a marketing claim into a time series. Publication of such a series is now a named verification event for the breadth-versus-depth forecast in Appendix~\ref{app:dashboard}; until it exists, every claim about ``ecosystem maturity''---including this book's---is graded [S].

\section{Democratizing AI through the LPDDR6 hardware arbitrage}
\label{sec:lpddr6}

The conventional Western paradigm assumes frontier training requires terabytes-per-second of High Bandwidth Memory---a technology where allied firms (SK~hynix, Samsung, Micron) hold a commanding lead and supply is structurally tight (HBM consumes roughly $3\times$ the wafer capacity of standard DRAM per bit~\cite{hbm-wafer}). US export controls accordingly treat HBM as a chokepoint. China's most interesting countermove is not to duplicate HBM---it is to sidestep it by exploiting massive capacities of low-cost commodity memory, accepting a bandwidth penalty and engineering around it in software. This is the LFP move (Chapter~\ref{ch:batteries}), executed in memory---and it is the training-centric variant of the Section~\ref{sec:template} template.

\subsection{The HBM scaling wall}

Analysis of large-scale Mixture-of-Experts training exposes the capacity limitation of the HBM-centric architecture. A feasibility study by the RIKEN Center for Computational Science of pretraining a 2.8-trillion-parameter MoE of the Kimi-K3 class finds that, under standard production-grade kernel fusion, sustaining a 20-PF(FP8) accelerator requires approximately \SI{14.5}{\tera\byte\per\second} of memory bandwidth per accelerator---beyond even HBM4 provisioning~\cite{riken-study}. The binding constraint, however, is capacity, not bandwidth: on a 128-accelerator cluster, the sharded model state alone (${\sim}\SI{306}{\giga\byte}$ per accelerator) exceeds the \SI{288}{\giga\byte} physical limit of current HBM4 stacks. The precise statement, with the discipline the result demands, is this: \emph{under the stated Kimi-K3-class model, optimizer representation, activation policy, and parallelization assumptions, the HBM4 configuration has a sub-256-accelerator capacity floor}---beneath it the Western architecture does not merely slow down but ceases to function~\cite{riken-study}. The threshold is configuration-specific, not universal: it moves with optimizer state, precision, sharding strategy, sequence length, expert placement, recomputation policy, and checkpoint design. It remains a strong result because the stated configuration is the \emph{standard} production configuration for the model class.

\subsection{CXMT and the LPDDR6 disruption}

ChangXin Memory Technologies (CXMT) supplies the memory class on which the alternative rests. Reports place its first LPDDR6 part---\SI{12800}{\mega\bit\per\second}, \SI{16}{\giga\bit} per die---at the tail end of R\&D verification in mid-2026, with sampling underway and mass production targeted for 2H26, which would place CXMT at or ahead of SK~hynix's LPDDR6 timeline: a Chinese memory firm reaching a leading-edge commodity node \emph{first}. This claim is graded [A]---analyst-reported roadmap evidence---because no official product announcement or customer qualification has been published; the falsification test is simply whether a qualified part ships on the reported schedule, and Appendix~\ref{app:dashboard} tracks it~\cite{cxmt-lpddr6}. The JEDEC SOCAMM2 module standard (LPDDR-on-module for AI servers, already in mass production for NVIDIA's own Grace/Vera-Rubin platforms; SK hynix began mass production of 192\,GB SOCAMM2 modules in April 2026~\cite{socamm2,skhynixSocamm}) provides the packaging: with approximately 18 SOCAMM2 modules, a single socket reaches ${\sim}\SI{2.3}{\tera\byte}$ of directly attached capacity---eight times the HBM4 ceiling, at a fraction of the cost per bit~\cite{riken-study}.

A second discipline separates what is \emph{observed} from what is \emph{proposed}. What the evidence supports: a capacity-first LPDDR architecture is a plausible and rational Chinese countermove, because it maps onto exactly the memory class where domestic capability is strongest and sidesteps the class where it is weakest. What the evidence does not yet support: that this architecture is an \emph{implemented} national countermove, or a product program that CXMT is leading. The design in this section and Appendix~\ref{app:lpddr} is this book's own engineering analysis [M] of what the pieces permit---not a report of a Chinese program, and the text is written so the two cannot be confused.

\subsection{The trade-off: capacity over peak bandwidth}

Before the trade is stated, the design hierarchy has to be fixed---and it is this book's own reconciled topology, rather than any preference for the novel corner, that fixes it. Appendix~\ref{app:lpddr}'s design point shows the pure-LPDDR socket clearing its $\sim$5\,TB/s requirement only under streaming-dominant access; under random expert-fetch behavior it misses by 10--30\%. The engineering ranking is therefore: \emph{(1) the hybrid HBM+LPDDR configuration as the baseline}---most of the capacity advantage, a quarter of the pin burden, graceful degradation; \emph{(2) pure LPDDR as a higher-risk, sanctions-contingent research configuration}---the corner a fully export-controlled program is pushed into, and the configuration whose feasibility milestone (the 5\,TB/s knee test) should be funded first precisely because it is the one in doubt; \emph{(3) HBM-only as the incumbent performance baseline}. What follows analyzes the pure-LPDDR corner because it is the \emph{strategically novel} one, not because it is the recommended one---the sanctions geography, not the engineering, is what selects it.

The arbitrage rests on a calculated hardware--software trade, and the entire subsection is an author-modeled result [M] pending the experimental program of Appendix~\ref{app:lpddr}. With \SI{2.3}{\tera\byte} per socket, activation recomputation can be eliminated entirely, cutting both FLOP and memory-traffic requirements; under aggressive kernel fusion the per-accelerator bandwidth requirement drops to a knee of ${\sim}\SI{5.0}{\tera\byte\per\second}$. An LPDDR6X system provisioned at \SI{7.0}{\tera\byte\per\second} peak delivers that achieved bandwidth, yielding a Model FLOP Utilization ceiling of about 67\%, against 89\% for the HBM4 configuration~\cite{riken-study}. The engineering price of that 7\,TB/s socket is itself nontrivial---on the order of 4,400 active data pins at LPDDR6 rates, several times the disclosed bandwidth of NVIDIA's own LPDDR-based Vera CPU (1.2\,TB/s)~\cite{vera-cpu}; Appendix~\ref{app:lpddr} does the arithmetic and states the validation milestones. What the MFU sacrifice buys: cost-per-bit an order of magnitude better; capacity-per-socket that \emph{enables} small-cluster training at all; and supply from a domestic fab ramping toward 300,000+ wafers per month (Chapter~\ref{ch:semi}). For the stated model class and parallelism layout, below 256 accelerators HBM4 capacity cannot hold the sharded state at all---there, LPDDR6X-class capacity is not the cheaper option but the enabling one.

\begin{figure}[htbp]
\centering
\begin{tikzpicture}
\begin{axis}[bookaxis, height=5.8cm,
  xmode=log, ymode=log,
  xlabel={Memory capacity per socket (GB, log)},
  ylabel={Memory bandwidth (TB/s, log)},
  title={The design space: capacity against bandwidth},
  xmin=20, xmax=9000, ymin=0.7, ymax=40,
  xtick={32,128,512,2048}, xticklabels={32,128,512,2048},
  ytick={1,2,4,8,16}, yticklabels={1,2,4,8,16},
  legend style={at={(0.02,0.03)}, anchor=south west},
]
\addplot[only marks, mark=*, mark size=2.6pt, color=sOne] coordinates {(192,8) (288,14.5)};
\addlegendentry{HBM-class accelerator}
\addplot[only marks, mark=square*, mark size=2.6pt, color=sTwo] coordinates {(2300,7) (1500,1.2)};
\addlegendentry{LPDDR-class socket}
\addplot[only marks, mark=triangle*, mark size=3.2pt, color=sFour] coordinates {(32,4)};
\addlegendentry{LX2 (accelerated CPU)}
\node[font=\scriptsize\color{inkSec}, anchor=east] at (axis cs:170,8) {B200};
\node[font=\scriptsize\color{inkSec}, anchor=west] at (axis cs:340,14.5) {HBM4 ceiling (288\,GB)};
\node[font=\scriptsize\color{inkSec}, anchor=south] at (axis cs:2300,7.8) {LPDDR6X proposal [M]};
\node[font=\scriptsize\color{inkSec}, anchor=south] at (axis cs:1500,1.35) {Vera (shipped)};
\node[font=\scriptsize\color{inkSec}, anchor=west] at (axis cs:36,3.6) {LX2};
\draw[dashed, color=inkMut, line width=0.6pt] (axis cs:306,0.7) -- (axis cs:306,20);
\node[font=\scriptsize\color{inkSec}, anchor=south east, align=right]
  at (axis cs:290,21) {model state per accelerator,\\128-node run ($\sim$306\,GB)};
\end{axis}
\end{tikzpicture}

\vspace{1mm}

\begin{tikzpicture}
\begin{axis}[bookaxis, height=4.6cm,
  xlabel={Provisioned memory bandwidth (TB/s)},
  ylabel={MFU ceiling (\%)},
  title={The capacity-for-bandwidth trade},
  xmin=1, xmax=16, ymin=0,
  xtick={2,5,7,10,14},
  legend style={at={(0.98,0.06)}, anchor=south east},
  ymax=112,
]
\addplot[color=sOne, mark=none, smooth] coordinates {
  (1,14) (2,28) (3,41) (4,53) (5,63) (6,66) (7,67) (9,72) (11,80) (14,89) (16,89)};
\addlegendentry{achievable MFU under aggressive fusion [M]}
\addplot[only marks, mark=*, mark size=2.6pt, color=sTwo] coordinates {(7,67)};
\addplot[only marks, mark=*, mark size=2.6pt, color=sFour] coordinates {(14,89)};
\node[font=\scriptsize\color{inkSec}, anchor=north west] at (axis cs:7.2,64) {LPDDR6X: 67\%};
\node[font=\scriptsize\color{inkSec}, anchor=north east] at (axis cs:13.8,86) {HBM4: 89\%};
\draw[dashed, color=inkMut, line width=0.5pt] (axis cs:5,0) -- (axis cs:5,63);
\node[font=\scriptsize\color{inkSec}, anchor=south west] at (axis cs:5.15,66) {knee $\approx$5\,TB/s};
\end{axis}
\end{tikzpicture}
\caption[Memory design space and the MFU trade]{Why capacity is the arbitrage. \emph{Above}: capacity against
bandwidth for shipped parts [V] and for the modeled LPDDR6X socket [M]. The dashed line is the per-accelerator
model state for the reference 2.8T-parameter sparse-MoE run at 128 nodes ($\sim$306\,GB): everything to its left
cannot hold the model at all, which is the sense in which the HBM path has a cluster-size floor. \emph{Below}:
the trade the substitution buys---22 points of MFU ceiling surrendered for an eightfold capacity gain and an
order-of-magnitude better cost per bit. Curve shape is modeled, not measured; Appendix~\ref{app:lpddr} states the
experiments that would confirm or refute it.}
\label{fig:memory}
\end{figure}


\subsection{Strategic consequence: decentralized pretraining}

The capacity advantage attacks the assumption that frontier pretraining requires centralized, gigawatt-class superclusters. Scale honesty first, because the intuitive figure for a machine of this class errs badly in the conservative direction. A 128-socket machine at 20\,PF(FP8) per socket is a 2.56-EF FP8 installation. Its power is \emph{not} tens of megawatts:

\[
P^{\mathrm{el}}_{\text{facility}} \;=\; \underbrace{N_{\text{sockets}}P^{\mathrm{el}}_{\text{socket}} + P^{\mathrm{el}}_{\text{network}} + P^{\mathrm{el}}_{\text{storage}}}_{\text{IT load}} \;+\; \underbrace{P^{\mathrm{el}}_{\text{cooling}} + P^{\mathrm{el}}_{\text{overhead}}}_{\text{facility overhead, i.e.\ }(\mathrm{PUE}-1)\times\text{IT load}} ,
\]
where $N_{\text{sockets}}$ is the number of accelerated processors, $P^{\mathrm{el}}_{\text{socket}}$ the power of one processor plus its attached memory tier, $P^{\mathrm{el}}_{\text{network}}$ and $P^{\mathrm{el}}_{\text{storage}}$ the interconnect and storage draw, and the last two terms the cooling and electrical losses that power-usage effectiveness (PUE) summarizes as a multiplier on IT load. Table~\ref{tab:128power} groups the terms exactly this way.

\begin{table}[htbp]
  \centering
  \small
  \caption{Power accounting for a 128-socket capacity-first pretraining cluster [M]. A ``socket'' is one accelerated processor plus its attached LPDDR6X memory tier; nodes carry two sockets.}
  \label{tab:128power}
  \setlength{\tabcolsep}{4pt}
  \begin{tabular}{lrl}
    \toprule
    Term & Value & Basis \\
    \midrule
    Compute silicon per socket & 2.0--2.6\,kW & $\sim$2$\times$ a B200-class part at 20\,PF(FP8) \\
    LPDDR6X tier per socket (2.3\,TB @ 7\,TB/s) & 0.3--0.5\,kW & $\sim$5--8\,pJ/bit at full stream \\
    \textbf{Socket total} & \textbf{2.3--3.1\,kW} & \\
    \midrule
    128 sockets (64 dual-socket nodes) & 295--397\,kW & \\
    Network (multi-rail, 1.6\,TB/s per socket) & 25--40\,kW & optics plus switching \\
    Storage and head nodes & 15--30\,kW & checkpoint tier at full state \\
    \textbf{IT load} & \textbf{335--467\,kW} & \\
    Cooling and facility overhead (PUE 1.15--1.30) & +50--140\,kW & liquid-cooled assumption \\
    \midrule
    \textbf{Facility total} & \textbf{$\sim$0.4--0.6\,MW} & 8--12 racks \\
    \bottomrule
  \end{tabular}
\end{table}

Half a megawatt is an ordinary industrial substation feeder, not a campus---which strengthens the decentralization argument rather than weakening it, and is the arithmetic a skeptical reader should check first. Even a 256-socket machine lands near 1\,MW. What that buys is \emph{smaller sovereign-scale}: a national-laboratory or mid-cap-company object; ``democratization'' here means the pretraining club expanding from three hyperscalers to every state and serious institution that wants in, and the capacity threshold itself is workload-specific, not a universal wall. Within that scope the claim stands [M]: sovereign states, national labs, and mid-size companies can execute full pretraining of frontier-class sparse models on 128--256-socket (64--128-node) clusters, because LPDDR6 provides the state capacity and Chinese sparse-model recipes (Chapter~\ref{ch:model}) provide the bandwidth frugality. Note the co-design loop closing at every scale: the same capacity-first, bandwidth-modest architecture that the \$2,000 enthusiast rig discovered bottom-up (Section~\ref{sec:quant}) reappears here as national strategy---Chinese models are trained sparse and bandwidth-lean \emph{because} of the hardware environment, and Chinese hardware is provisioned capacity-rich \emph{because} the models reward it. The Western trillion-dollar bet on centralized HBM-bound data centers is thus exposed on two flanks at once: its scarce input (HBM) is bypassed, and its scale premise (only hyperscalers can pretrain) is dissolved. Hardware proliferation plus software commoditization equals the democratization---on Chinese terms---of the means of AI production.

\section{The West defects to the open ISA: NVIDIA, Google, Meta}
\label{sec:west-riscv}

The final structural fact of the compute element is that the Western ecosystem is not defending the proprietary-ISA order---it is quietly joining the open one. The evidence spans every layer of the US stack.

A taxonomy first, because ``independence'' has five distinct layers and conflating them flatters the thesis: (1) host-ISA independence---what RISC-V directly delivers; (2) accelerator-ISA and microarchitecture independence; (3) framework and compiler portability; (4) scale-up/scale-out fabric independence; (5) memory and packaging independence. RISC-V settles only layer~1; the larger AI rents live in layers 2--5, which is where CANN/UB open-sourcing, the matrix-extension standardization, and the memory strategy of this chapter respectively operate. Unit forecasts deserve the same discipline: most RISC-V units are embedded controllers, and unit share is not server-revenue share---though the SHD revenue decomposition ($\sim$70\% of projected 2031 RISC-V \emph{revenue} from AI-accelerator SoCs) shows the value migrating to exactly the contested layer~\cite{shd-2026}.

\textbf{NVIDIA} is the most striking case. Its GPUs have carried embedded RISC-V management cores since 2016---ten to forty per chip, roughly one billion cores shipped in 2024 alone~\cite{nvidia-riscv-cores}. In July 2025, at the RISC-V Summit \emph{China} in Shanghai, NVIDIA announced that CUDA itself is being ported to RISC-V host CPUs, making the open ISA the third sanctioned CUDA host after x86 and Arm---explicitly enabling ``a RISC-V CPU to be the main application processor'' of a CUDA system~\cite{cuda-riscv}. Two readings of this move are available and both should be on the record. The moat-extension reading: NVIDIA is porting CUDA \emph{onto} the open host layer precisely to keep the accelerator rent (layers 2--3) intact---the RISC-V CPU runs Linux and drivers while the GPU keeps the compute~\cite{cuda-riscv-analysis}. The concession reading: the platform owner no longer expects the proprietary host order to hold, and is repositioning its rent one layer up---at a summit in Shanghai, weeks after its China market went to zero. The readings agree on the structural fact this book needs: the foundation layer is going open with both superpowers' participation, and the contest has moved to the layers above it. In January 2026 SiFive announced NVLink Fusion integration for datacenter-class RISC-V subsystems~\cite{sifive-nvlink}; RISC-V International's own five-year CEO, Calista Redmond, departed in December 2024---for NVIDIA~\cite{redmond}.

\textbf{Google} co-founded the RISC-V Foundation in 2015 and the RISE software consortium in 2023, ships RISC-V in its Titan security silicon, deploys SiFive X280 vector cores inside its TPU ecosystem as the programmable front-end to its matrix units, and committed Android to the architecture~\cite{google-riscv}. \textbf{Meta}'s MTIA accelerator family---four generations disclosed through March 2026, scaling to 10\,PFLOPS parts deployed 72 to a rack---is built on RISC-V processing elements throughout~\cite{mtia}. \textbf{Microsoft} is the conspicuous holdout, its Azure silicon remaining Arm-based (Cobalt) with no public RISC-V commitment---a gap worth watching precisely because it is now singular~\cite{cobalt}. Attempts in Washington to restrict RISC-V collaboration---congressional letters in 2023, 2024, and 2025 noting that Chinese firms hold nearly half the consortium's board seats---have produced no enforceable action, for the reason the Commerce Department itself conceded: the standard is open, Swiss-domiciled, and restricting US participation would harm US firms without slowing China at all~\cite{riscv-restrict}.

The pattern matters more than any single datum: the ISA layer is tipping the way the model layer tipped. Forecasts now put RISC-V at $\sim$25\% of processor units by 2030, with AI as the driver~\cite{shd-2026,omdia}. When the platform owner of the AI era ports its crown-jewel software to the open ISA at a summit in Shanghai, the standards war is not being lost; it is being conceded. What remains proprietary in the Western stack---the GPU microarchitecture, CUDA's implementation, Arm's cores---sits on top of a foundation both superpowers now agree will be open. Only one of them has organized its national industry to own that foundation.

\section{Assessment}

Honesty about the current balance, and about what this chapter does \emph{not} claim (Section~\ref{sec:notobsolete}): NVIDIA-plus-HBM remains the superior machine per watt for dense frontier training today (CloudMatrix-class systems reach parity by brute force at $\sim$4$\times$ the power~\cite{cloudmatrix}); Ascend software maturity cost DeepSeek a flagship cycle~\cite{ascend-delay}; domestic accelerator output in 2026 (roughly 1.5--2M die-equivalents) remains an order of magnitude below NVIDIA's~\cite{semianalysis-ascend,ifp-b30a}; and the LX2's AI-serving numbers are, as its own analysts insist, projections from published specifications pending direct measurement~\cite{dhm-long}. But the structural position has inverted since this book's first draft. A GPU-free Chinese machine leads both TOP500 benchmarks; the CPU-versus-accelerator distinction has been shown, quantitatively, to be an allocation choice rather than a category; the template is portable to an ISA no one can revoke, in two memory variants matched to the two halves of the AI workload; and the Western ecosystem itself---NVIDIA first among them---is porting its software to that ISA. The remaining dependency of the compute element is physical: wafers, HBM stacks, and DRAM bits. That is the third element.
